\section{Introducció}

\subsection{Context}

Aquest projecte s'ha desenvolupat en el marc del \textit{hackathon} de l'assignatura de Sistemes Multimèdia del Grau en Sistemes Multimèdia de la Universitat Autònoma de Barcelona, durant el curs 2025--2026. L'objectiu del hackathon és dissenyar i implementar un sistema complet que integri múltiples tecnologies multimèdia i de computació al núvol en un termini limitat.

El projecte \textbf{Dal-i} —nom que fa referència a l'artista Salvador Dalí i al model d'IA generativa DALL-E— és un sistema que permet a qualsevol usuari crear dibuixos artístics a través d'un braç robòtic, sense necessitat de cap habilitat artística prèvia. El sistema accepta una imatge d'entrada i un estil artístic (introduïts per veu, text o selecció), processa la imatge mitjançant tècniques de visió per computador, i genera les coordenades de trajectòria que el robot ha de seguir per reproduir el dibuix.

\subsection{Motivació}

La motivació principal del projecte és construir un pont entre tres àmbits tecnològics que rarament es combinen en un sol sistema educatiu:

\begin{itemize}
  \item \textbf{Intel·ligència artificial generativa}: ús de models de difusió i transformadors per generar o transferir estils artístics.
  \item \textbf{Visió per computador clàssica}: detecció de contorns, extracció de característiques i optimització de trajectòries per a robòtica.
  \item \textbf{Arquitectura cloud-native}: desplegament escalable i de baix cost mitjançant Google Cloud Platform, amb separació de responsabilitats en microserveis.
\end{itemize}

A més, l'addició de suport multilingüe (10 idiomes) i d'entrada per veu reflecteix l'interès del grup per fer el sistema accessible a usuaris de qualsevol llengua i entorn.

\subsection{Objectius}

Els objectius principals del projecte són:

\begin{enumerate}
  \item Implementar un \textit{pipeline} complet de processament d'imatges que converteixi qualsevol fotografia en coordenades de trajectòria per a un braç robòtic SCARA.
  \item Desplegar el sistema sobre una arquitectura \textit{cloud-native} escalable, fent ús de Cloud Run i Cloud Functions de Google Cloud Platform.
  \item Oferir múltiples modalitats d'entrada: selecció d'estil per text, per veu i per botons de selecció ràpida.
  \item Integrar IA generativa (Gemini 2.5 Flash) per a la transferència d'estil visual en un mode avançat opcional.
  \item Donar suport multilingüe automàtic per permetre l'ús del sistema en 10 idiomes.
  \item Proporcionar una experiència d'usuari fluida amb historial persistent per sessió, paginació i esborrat en cascada.
\end{enumerate}

\subsection{Estructura de l'informe}

La resta d'aquest document s'organitza de la manera següent. La secció~\ref{sec:descripcio} descriu l'aplicació des del punt de vista de l'usuari. La secció~\ref{sec:arquitectura} detalla l'arquitectura del sistema. La secció~\ref{sec:desenvolupament} explica el procés de desenvolupament per fases. La secció~\ref{sec:tecnologies} llista les tecnologies utilitzades. La secció~\ref{sec:problemes} recull els problemes trobats i les solucions adoptades. La secció~\ref{sec:resultats} presenta els resultats obtinguts. Finalment, la secció~\ref{sec:conclusions} recull les conclusions i el treball futur.
